\documentclass[../../main2.tex]{subfiles}

\begin{document}

	\chapter*{引言}
	\addcontentsline{toc}{chapter}{引言}

	\begin{motivation}
		本书的前身《社会动力学与公理化体系》回答了「为什么」：经济学讲激励，政治学讲权力，人类学讲意义——它证明了这些能被放进同一条回溯的解释链。\\
		但它回答不了「接下来呢」。\\
		\textbf{核心动机}：把问法换掉——不再问「为什么是这个结果」，改问「这个结果里，各因素各占多少」。前者是解释学，后者是预测学；这一版要建立的是后者。
	\end{motivation}

	\begin{logicmap}
	\begin{tikzpicture}[node distance=0.9cm, every node/.style={draw, rectangle, rounded corners, align=center}]
	\node (q) {写作动机：从「为什么」到「接下来呢」};
	\node (r) [below=of q] {面向读者群体};
	\node (s) [below=of r] {新结构：一条前向的链};
	\node (h) [below=of s] {如何阅读本书};
	\node (c) [below=of h] {检验标准：一句话立论};
	\draw[-{Stealth}] (q) -- (r);
	\draw[-{Stealth}] (r) -- (s);
	\draw[-{Stealth}] (s) -- (h);
	\draw[-{Stealth}] (h) -- (c);
	\end{tikzpicture}
	\end{logicmap}

	\section{写作动机：从「为什么」到「接下来呢」}

	\begin{readingnote}
		你有没有想过，为什么历史书能告诉你「罗马为什么亡」，却没有一本历史书能告诉你「谁正在重演罗马」？\\
		解释过去和预测未来，用的是同一堆事实——为什么一个总能做到，另一个总是失败？
	\end{readingnote}

	一切始于一连串深夜的讨论。朋友们的质疑像种子一样落地生根：社会科学能不能像物理一样，从一组基本原理出发推导出结论？前身给出了肯定的答案——从地理的初始条件，推到结构，再推到经济、文化、政治与法律，一条完整的回溯链。

	但作者在自审报告里亲手承认了这条链的限度：它解释了「为什么是这样」，解释不了「接下来会怎样」。\textbf{解释和预测是不对称的}——解释可以事后挑原因，预测必须事前排权重。解释罗马亡于军团膨胀，不需要说出「下一个是谁」；要说出「下一个是谁」，就必须回答「军团膨胀占几成、财政占几成、瘟疫占几成」。

	所以这一版的写作动机，不再是展示「如何问为什么」，而是展示「如何问各占多少」。这个问法的转换是全书性的：从\textbf{回溯}换成\textbf{前向}，从「为什么是这个结果」换成「这个结果里，各因素各占多少」。只有能定各占多少，才能回答「如果当初没有 X，今天会怎样」——而反事实推演（counterfactual），正是预测的最小单元。

	\paragraph*{逻辑衔接}
	我们已经明确了新的目标：预测。但一本讲预测的书，首先得说清自己为谁而写、需要什么背景——否则「预测」会被误读成算命。下一节「面向读者群体」处理这个问题。

	\section{面向读者群体}

	\begin{readingnote}
		你有没有想过，为什么同一本书，有人觉得太数学看不懂，有人觉得不够数学没深度？\\
		这反映的其实是\textbf{读者背景的多样性}。一本讲「预测的社会科学」的书，怎么才能让不同背景的人都觉得有收获？
	\end{readingnote}

	本书仍然在两个极端之间寻找平衡：数理严谨性与社科可读性。这次的天平更难摆——因为「预测」天然要求数字，而「社会」天然拒绝数字。本书的处理是：正文用自然语言推进直觉，每个关键概念配一个「数学建模」盒子放形式化——可跳过，但跳过之后你要接受「知其然」的折扣。

	因此，本书的目标读者包括：
	\begin{itemize}
		\item 希望理解社会系统如何被\textbf{前向推导}的自然科学背景者
		\item 愿意接受概率工具、但不希望被公式淹没的社会科学研究者
		\item 对「社会能不能预测」这个问题本身感兴趣的普通读者
	\end{itemize}

	第三类读者请特别注意：本书的「预测」是\textbf{概率断言}，不是\textbf{预言}。这个区别将在第6章被严格展开——这里先埋一句：说「降水概率 70\%」的气象学家，和说「明天必然下雨」的巫师，做的是两种完全不同的事。

	\paragraph*{逻辑衔接}
	我们已经明确了本书为谁而写、以及「预测」这个词的用法边界。接下来要回答的是结构问题：这条前向的链长什么样？各部分之间是什么依赖关系？下一节「新结构」给出全书地图。

	\section{新结构：一条前向的链}

	\begin{readingnote}
		你有没有想过，为什么有的书读完只记得零散的观点，有的书读完却带走一张地图？\\
		差别往往不在内容，而在\textbf{结构}——章节之间是并列摆放，还是彼此咬合？
	\end{readingnote}

	前身的一个教训（作者自审确认）是：四个领域被写成了并列的四章，像四本教科书的合订本。新结构对这一点做了根治：\textbf{全书是一条链，每一环都是下一环的前提}。

	\begin{figure}[h]
	\centering
	\begin{tikzpicture}[
		node distance=0.7cm and 1.6cm,
		every node/.style={draw, rectangle, rounded corners, align=center, font=\small}
	]
	\node (p1) {Part I\\动机与目标};
	\node (p2) [below=of p1] {Part II\\演化状态空间};
	\node (p3) [below=of p2] {Part III\\条件概率与贡献};
	\node (p4) [below=of p3] {Part IV\\四个通道};
	\node (p5) [below=of p4] {Part V\\预测与循环};
	\draw[-{Stealth}] (p1) -- (p2);
	\draw[-{Stealth}, very thick] (p2) -- node[right, draw=none]{硬依赖} (p3);
	\draw[-{Stealth}] (p3) -- (p4);
	\draw[-{Stealth}] (p4) -- (p5);
	\draw[-{Stealth}, dashed] (p5.east) .. controls +(2.2,0) and +(2.2,0) .. node[right, draw=none, align=left]{循环：预测\\塑造新动机} (p1.east);
	\end{tikzpicture}
	\caption{全书五部依赖图。Part II $\rightarrow$ Part III 是硬依赖：条件概率必须定义在状态空间上——这正是原书附录断裂点的修复。}
	\label{fig:intro-dag}
	\end{figure}

	这条链的每一站，都长在上一站遗留的问题上：
	\begin{itemize}
		\item \textbf{动机与目标}（Part I）：任何解释的最小起点是什么？动机为什么先于结构？动机不决定结果，那中间隔着什么？——隔着「目标」。
		\item \textbf{演化状态空间}（Part II）：要预测未来，先得规定「未来由什么构成」——物质、人、相互作用。这是原书缺失的\textbf{本体论}层。
		\item \textbf{条件概率与贡献分解}（Part III）：未来不可唯一确定时，如何定量说话？$C_A = P(D \mid A) - P(D \mid \neg A)$——「各占多少」在这里获得定义。
		\item \textbf{四个通道}（Part IV）：抽象的「贡献」分解成具体通道——模因、经济、政治、法律。\textbf{不是并列，是嵌套}：每章是上一章「加入一个约束」。
		\item \textbf{预测与循环}（Part V）：把贡献表拼回时间线，得到一条可辩护的未来——然后，预测出的未来塑造下一轮的动机。循环闭合。
	\end{itemize}

	写作上有两条铁律贯穿全书，值得在这里先打招呼。其一，\textbf{每章开头都是上一章遗留的问题}——你不会看到「本章讨论 X」这样的开头。其二，\textbf{每个新概念都通过「加入一个约束」导出}——如果 A 不需要 X，性质是 P；加入约束 X，所有性质随之改变，这就是 B 从 A 中分离出来的时刻。

	\paragraph*{逻辑衔接}
	结构图有了。但一条环环相扣的链，会不会把读者锁死在单一读法里？前身已经论证过「正向建构与反向学习不冲突」——这里只需交代本书的具体读法。下一节「如何阅读本书」给出阅读策略。

	\section{如何阅读本书}

	\begin{readingnote}
		你有没有遇到过这样的情况：有些书\textbf{结构清晰}但读起来枯燥，有些书\textbf{生动有趣}却杂乱无章？\\
		如果一本书想既保持逻辑统一、又允许灵活阅读——这可能吗？
	\end{readingnote}

	\subsection{线性阅读路径}

	如果这是你第一次接触这类前向推导的思考，建议按目录顺序读。这条链的依赖关系是真实的：不先有状态空间（Part II），条件概率（Part III）就是悬空的符号——原书的断裂正是出在这里，本书的顺序就是修复方案本身。

	\subsection{跳跃阅读路径}

	如果你已有相关背景，可以按问题跳读。每一章开头的「逻辑地图」标明该章位置与前后依赖；每一节末尾的「逻辑衔接」显式写出「下一节将解决什么」——它们合起来构成一张可以脱离正文使用的全书地图。想直接看「各占多少」的读者，可以从第7章进入，缺什么再回第3章补。

	\subsection{正向建构与反向学习}

	从前身继承的结论在这里依然成立：从基本原理向下读（正向建构）和从具体问题向上回溯（反向学习），是\textbf{对同一结构的两种线性遍历方式}，并不冲突。唯一的新增提醒是：本书的硬依赖边（Part II $\rightarrow$ Part III）在反向阅读时最容易踩空——从第7章往回跳时，请务必在第3章停留，确认「状态」和「相互作用」的定义在手。

	\paragraph*{逻辑衔接}
	读法交代完毕。还剩最后一件事：把「这本书成功了」说清楚——否则「预测」会滑向它最俗的意义。下一节给出全书唯一的检验标准。

	\section{检验标准：一句话立论}

	\begin{readingnote}
		你有没有想过，为什么听了一场「大势分析」之后热血沸腾，三个月后发现什么也没预测对？\\
		因为\textbf{没有可检验的断言}的预测，只是情绪按摩。一本书的预测力，必须有验收标准。
	\end{readingnote}

	本书的验收标准只有一句话：

	\begin{quote}
		\textbf{如果一个陌生人读完这本书，他能用书里的公式，对一个他没见过的社会事件给出概率判断，并说出「这个判断依赖于哪几个因素」——那么这本书就成功了。}
	\end{quote}

	这句话取代了前身的检验标准。前身的标准只验「统一视角是否说清」；新标准必须验「预测能力是否建立」。两者的差别就是本书存在的理由。

	相应地，「我们能说什么 / 我们不能说什么」必须分开写。本书能给的：对一个事件，一组因素的贡献排序与概率断言；对一个「如果当初没有 X」的问题，一条可辩护的反事实推演。本书不能给的：明天的股价、下一位候选人、战争的具体日期。第13章的失败模式清单会把这个「不能」列满——少于三条的失败清单不值得信任。

	\paragraph*{逻辑衔接}
	标准立好了。接下来正式进入正文——但在出发前，先看一眼行军路线：各章要解决什么问题、彼此怎么咬合。下一节「章节预告」给出全书路标，然后我们从最日常的一句话开始：「他这么做，是因为……」。

	\section{章节预告}

	\begin{itemize}
		\item \hyperref[ch:motivation]{第1章 动机：解释的最小起点}——「他这么做是因为……」这句话里隐藏了什么假设？任何解释的最小单元是主体、动机、环境；动机先于结构；但动机不决定结果。
		\item \hyperref[ch:goal]{第2章 目标：把意向写成可比较的量}——「想要 X」里的 X 是什么形状的东西？排序、不可通约性、效用函数、显示性偏好（revealed preference），最后退无可退：交给概率。
		\item \hyperref[ch:state]{第3章 状态三元组}——「未来」这个词，说的是什么在变？物质、人、相互作用，以及为什么三个就够了。
		\item \hyperref[ch:survival]{第4章 人存}——所有目标里，有没有一个是不能放弃的？人存作为一切目标函数的公共约束域。
		\item \hyperref[ch:evolution]{第5章 演化}——状态怎么动？有压扩增与无压扩增的分叉，以及可预测窗口。
		\item \hyperref[ch:prediction]{第6章 预测的数学形式}——「预测」在数学上到底是什么？条件概率链的乘积分解。
		\item \hyperref[ch:contribution]{第7章 贡献分解}——A 和 B 共同导致了 D，各占多少？$C_A = P(D \mid A) - P(D \mid \neg A)$。
		\item \hyperref[ch:channels]{第8章 通道生成器}——所有社会行为，按「复制还是转移」分只有两类。
		\item \hyperref[ch:meme]{第9章 模因}、\hyperref[ch:economy]{第10章 经济}、\hyperref[ch:politics]{第11章 政治}、\hyperref[ch:law]{第12章 法律}——四个通道，逐层加约束：稀缺约束、身份权重、定责编码。\textbf{嵌套，不是并列}。
		\item \hyperref[ch:forecast]{第13章 从贡献表到未来走向}——怎么把分解结果拼回一条时间线？包含全书唯一一次完整演示。
		\item \hyperref[ch:closure]{第14章 循环的闭合}——预测出的未来，会不会反过来塑造新的动机？
	\end{itemize}

	\paragraph*{逻辑衔接}
	路线图已经摊开。现在，让我们从一个最日常的动作开始——解释。当你说出「他这么做，是因为……」的那一刻，你已经站在了整本书的起点上。第1章「动机：解释的最小起点」将拆开这句话，看看里面藏着什么假设。

\end{document}
